\begin{abstract}
LLM judges are the de facto standard for evaluating agentic tool-calling systems, yet their reliability in structured, dependency-driven workflows remains largely unexamined. We present AgentJudgeBench, the first benchmark to systematically study LLM-as-a-judge behavior over workflow DAGs. The benchmark comprises 3,808 instances spanning six DAG topologies and three difficulty tiers, evaluated across five generators (3B–70B open-weight models and GPT-5.4), six judges (20B to frontier scale), and both with- and no-ground-truth settings.
Our analysis uncovers a set of consistent and previously unreported phenomena. First, judge alignment degrades monotonically with task difficulty, with degradation occurring 1.5× faster in the absence of ground truth. Second, in no-GT settings, alignment converges to a narrow 77–82\% band regardless of judge scale, revealing a structural ceiling that cannot be overcome by model capacity alone. Third, providing ground truth can hurt frontier judges: GPT-5.4 drops by 1.5 pp and Gemini-2.5-Pro by 3.9 pp, consistent with over-anchoring effects.
Finally, we evaluate common mitigation strategies. Chain-of-thought reasoning yields negligible gains (less than 0.3 pp), while structured evaluation rubrics improve alignment by +4.8–+6.5 pp over free-form prompting. With ground truth, QwQ-32B achieves the strongest performance; without it, frontier models lead only marginally within the no-GT ceiling.
Together, our results identify fundamental limitations of current LLM judges and establish practical guidelines for reliable evaluation in agentic systems.
\end{abstract}